% Inline TikZ versions of the policy-architecture figures used by
% invman_lostsales.tex. These macros intentionally avoid pre-rendered PDFs.

\colorlet{invblue}{blue!70!black}
\colorlet{invgreen}{green!45!black}
\colorlet{invorange}{orange!75!black}
\colorlet{invred}{red!75!black}
\colorlet{invgray}{black!55}

\tikzset{
  invstage/.style={
    rounded corners=3pt,
    thick,
    align=center,
    inner sep=5pt,
    minimum height=11mm
  },
  invinput/.style={invstage,draw=invgreen!70!black,fill=invgreen!10},
  invmiddle/.style={invstage,draw=invblue!70!black,fill=invblue!8},
  invdecoder/.style={invstage,draw=invorange!80!black,fill=invorange!10},
  invoutput/.style={invstage,draw=invred!70!black,fill=invred!10},
  invflow/.style={-{Latex[length=3.2mm,width=2.2mm]},thick,invblue},
  invsoft/.style={circle,draw=invblue!70!black,fill=invblue!10,thick,
                  minimum size=9mm,inner sep=0pt},
  invleaf/.style={invstage,draw=invgreen!65!black,fill=invgreen!8,
                  minimum width=19mm},
  invbox/.style={draw=invblue!55,dashed,rounded corners=4pt,inner sep=8pt},
  invevsource/.style={invstage,draw=invgreen!70!black,fill=invgreen!10,
                      minimum width=27mm,minimum height=13mm,font=\small},
  invevstock/.style={invstage,draw=invblue!70!black,fill=invblue!8,
                     minimum width=29mm,minimum height=13mm,font=\small},
  invevdemand/.style={invstage,draw=invorange!80!black,fill=invorange!10,
                      minimum width=28mm,minimum height=13mm,font=\small},
  invevlost/.style={invstage,draw=invred!70!black,fill=invred!8,
                    minimum width=28mm,minimum height=13mm,font=\small},
  invmaterial/.style={-{Latex[length=3.4mm,width=2.4mm]},very thick,invblue},
  invrequest/.style={-{Latex[length=3.1mm,width=2.2mm]},semithick,dashed,invgreen!70!black},
  invdemandflow/.style={-{Latex[length=3.4mm,width=2.4mm]},very thick,invorange!85!black},
  invlostflow/.style={-{Latex[length=3.4mm,width=2.4mm]},very thick,dashed,invred},
  invlabel/.style={font=\footnotesize,align=center,fill=white,inner sep=1.5pt},
  invcost/.style={font=\footnotesize,align=center}
}

\newcommand{\PolicyPipelineFigure}{%
\resizebox{0.96\linewidth}{!}{%
\begin{tikzpicture}[
  x=1cm,y=1cm,
  stagebox/.style={minimum width=28mm,minimum height=14mm,font=\small}
]
  \node[invinput,stagebox] (state) at (0,0)
    {\textbf{State}\\[1pt]$\mathbf S_t$};
  \node[invmiddle,stagebox] (norm) at (3.45,0)
    {\textbf{Normalize}\\[1pt]$\mathbf x=\mathbf S_t/\kappa$};
  \node[invmiddle,stagebox] (backbone) at (6.95,0)
    {\textbf{Backbone}\\[1pt]\footnotesize linear, NN, tree};
  \node[invdecoder,stagebox] (decoder) at (10.45,0)
    {\textbf{Decoder}\\[1pt]\footnotesize direct, ordinal, tree};
  \node[invoutput,stagebox] (action) at (13.9,0)
    {\textbf{Order}\\[1pt]$q_t$};

  \draw[invflow] (state) -- (norm);
  \draw[invflow] (norm) -- (backbone);
  \draw[invflow] (backbone) -- (decoder);
  \draw[invflow] (decoder) -- (action);

  \begin{scope}[on background layer]
    \node[invbox,fit=(norm)(backbone)(decoder),inner sep=7pt] (policy) {};
  \end{scope}
  \node[invblue,fill=white,inner sep=1.5pt,below=5pt of policy.south]
    {\footnotesize learned policy $\pi_{\btheta}$, optimized by CMA-ES};
\end{tikzpicture}}%
}

\newcommand{\LinearPolicyFigure}{%
\resizebox{0.78\linewidth}{!}{%
\begin{tikzpicture}[x=1cm,y=1cm]
  \node[invinput,minimum width=28mm,minimum height=17mm,font=\small] (x) at (0,0)
    {normalized state\\$\mathbf x$};
  \node[invmiddle,minimum width=34mm,minimum height=17mm,font=\small] (affine) at (4.0,0)
    {affine map\\$W\mathbf x+b$};
  \node[invoutput,minimum width=28mm,minimum height=17mm,font=\small] (u) at (8.0,0)
    {latent outputs\\$\mathbf u$};

  \draw[invflow] (x) -- (affine);
  \draw[invflow] (affine) -- (u);
  \node[invblue,below=4pt of affine.south,font=\footnotesize,align=center]
    {linear backbone: no hidden layer};
\end{tikzpicture}}%
}

\newcommand{\NeuralPolicyFigure}{%
\resizebox{0.88\linewidth}{!}{%
\begin{tikzpicture}[x=1cm,y=1cm]
  \node[invinput,circle,minimum size=8mm,inner sep=0pt,font=\footnotesize] (x1) at (0,1.1) {$x_1$};
  \node[invinput,circle,minimum size=8mm,inner sep=0pt,font=\footnotesize] (x2) at (0,0.25) {$x_2$};
  \node[font=\small,invgray] (xdots) at (0,-0.55) {$\vdots$};
  \node[invinput,circle,minimum size=8mm,inner sep=0pt,font=\footnotesize] (xL) at (0,-1.35) {$x_L$};

  \node[invmiddle,circle,minimum size=8mm,inner sep=0pt,font=\footnotesize] (h1) at (3.2,1.55) {$h_1$};
  \node[invmiddle,circle,minimum size=8mm,inner sep=0pt,font=\footnotesize] (h2) at (3.2,0.65) {$h_2$};
  \node[font=\small,invgray] (hdots) at (3.2,-0.2) {$\vdots$};
  \node[invmiddle,circle,minimum size=8mm,inner sep=0pt,font=\footnotesize] (hH) at (3.2,-1.15) {$h_8$};

  \node[invoutput,circle,minimum size=8mm,inner sep=0pt,font=\footnotesize] (u1) at (6.4,0.95) {$u_1$};
  \node[font=\small,invgray] (udots) at (6.4,0.05) {$\vdots$};
  \node[invoutput,circle,minimum size=8mm,inner sep=0pt,font=\footnotesize] (uM) at (6.4,-0.95) {$u_m$};

  \foreach \i in {x1,x2,xL}{
    \foreach \j in {h1,h2,hH}{
      \draw[invblue!35,thin] (\i) -- (\j);
    }
  }
  \foreach \i in {h1,h2,hH}{
    \foreach \j in {u1,uM}{
      \draw[invblue!35,thin] (\i) -- (\j);
    }
  }

  \node[invgreen,above=2pt of x1.north,font=\footnotesize] {input $\mathbf x$};
  \node[invblue,above=2pt of h1.north,font=\footnotesize] {hidden layer};
  \node[invred,above=2pt of u1.north,font=\footnotesize] {latent $\mathbf u$};
  \node[invblue,below=5pt of hH.south,font=\footnotesize,align=center]
    {$\mathbf u=V\phi(U\mathbf x+c)+e$, with $H=8$};
\end{tikzpicture}}%
}

\newcommand{\DecoderHeadsFigure}{%
\resizebox{\linewidth}{!}{%
\begin{tikzpicture}[x=1cm,y=1cm]
  \node[invmiddle,minimum width=24mm,minimum height=19mm] (latent) at (0,-2.7)
    {latent\\outputs\\$\mathbf u$};

  \node[invdecoder,anchor=west,text width=73mm] (direct) at (2.4,0)
    {\textbf{soft-gated direct}\\
     $q_t=\left\lfloor\sigma(g(\mathbf x))\,\mathrm{softplus}(r(\mathbf x))\right\rceil_{\ge0}$};
  \node[invdecoder,anchor=west,text width=73mm] (ordinal) at (2.4,-1.8)
    {\textbf{soft-gated ordinal}\\
     $q_t=\left\lfloor\sigma(g(\mathbf x))\sum_{k=1}^{M}\sigma(o_k(\mathbf x))\right\rceil_{\ge0}$};
  \node[invdecoder,anchor=west,text width=73mm] (tree) at (2.4,-3.6)
    {\textbf{soft tree with linear leaves}\\
     $q_t=\left\lfloor\sum_{\ell}\pi_{\ell}(\mathbf x)\,q_{\ell}(\mathbf x)\right\rceil_{\ge0}$};
  \node[invdecoder,anchor=west,text width=73mm] (categorical) at (2.4,-5.4)
    {\textbf{categorical baseline}\\
     $q_t=\argmax_{q\in\{0,\ldots,\bar q\}} y_q$};

  \node[invoutput,minimum width=24mm,minimum height=19mm] (order) at (12.9,-2.7)
    {\textbf{integer}\\\textbf{order}\\$q_t$};

  \foreach \head in {direct,ordinal,tree,categorical}{
    \draw[invflow] (latent.east) -- (\head.west);
    \draw[invflow] (\head.east) -- (order.west);
  }
  \node[invgray,below=3mm of categorical.south,align=center]
    {\footnotesize The reported new experiments use the direct, ordinal, and tree decoders; the categorical head is the earlier score-per-order representation.};
\end{tikzpicture}}%
}

\newcommand{\SoftTreeFigure}{%
\resizebox{0.82\linewidth}{!}{%
\begin{tikzpicture}[x=1cm,y=1cm]
  \node[invinput,minimum width=34mm] (x) at (0,2.0)
    {normalized state\\$\mathbf x=\mathbf S_t/\kappa$};
  \node[invsoft] (root) at (0,0) {$n_0$};
  \node[invsoft] (left) at (-2.4,-1.8) {$n_1$};
  \node[invsoft] (right) at (2.4,-1.8) {$n_2$};
  \node[invleaf] (l1) at (-3.6,-3.7) {$q_1(\mathbf x)$};
  \node[invleaf] (l2) at (-1.2,-3.7) {$q_2(\mathbf x)$};
  \node[invleaf] (l3) at (1.2,-3.7) {$q_3(\mathbf x)$};
  \node[invleaf] (l4) at (3.6,-3.7) {$q_4(\mathbf x)$};
  \node[invoutput,minimum width=45mm] (out) at (0,-5.8)
    {$q_t=\left\lfloor\sum_{\ell}\pi_{\ell}(\mathbf x)q_{\ell}(\mathbf x)\right\rceil_{\ge0}$};

  \draw[invflow] (x) -- (root);
  \draw[invflow] (root) -- node[above left,font=\footnotesize] {$1-\sigma(\cdot)$} (left);
  \draw[invflow] (root) -- node[above right,font=\footnotesize] {$\sigma(\cdot)$} (right);
  \draw[invflow] (left) -- (l1);
  \draw[invflow] (left) -- (l2);
  \draw[invflow] (right) -- (l3);
  \draw[invflow] (right) -- (l4);

  \begin{scope}[on background layer]
    \node[invbox,fit=(l1)(l2)(l3)(l4)] (leaves) {};
  \end{scope}
  \draw[invflow] (leaves.south) -- (out.north);
  \node[invgray,below=2pt of out.south,align=center]
    {\footnotesize Path probabilities $\pi_{\ell}(\mathbf x)$ weight the leaf quantities before rounding.};
\end{tikzpicture}}%
}

\newcommand{\CMAESLoopFigure}{%
\resizebox{\linewidth}{!}{%
\begin{tikzpicture}[x=1cm,y=1cm]
  \node[invmiddle,minimum width=31mm] (dist) at (0,0)
    {search distribution\\$\mathcal N(\bmu_i,\sigma_i^2\bSigma_i)$};
  \node[invinput,minimum width=30mm] (sample) at (3.9,0)
    {sample\\$\btheta_i^{(1)},\ldots,\btheta_i^{(N)}$};
  \node[invdecoder,minimum width=34mm] (eval) at (7.9,0)
    {simulate policies\\estimate $\hat C_T(\btheta)$};
  \node[invmiddle,minimum width=26mm] (rank) at (11.5,0)
    {rank\\by cost};
  \node[invoutput,minimum width=31mm] (update) at (14.8,0)
    {update\\$\bmu,\sigma,\bSigma$};

  \draw[invflow] (dist) -- (sample);
  \draw[invflow] (sample) -- (eval);
  \draw[invflow] (eval) -- (rank);
  \draw[invflow] (rank) -- (update);
  \draw[invflow] (update.south) -- ++(0,-1.1) -| (dist.south);

  \node[invorange,above=2pt of eval.north] {\footnotesize simulator call};
  \node[invblue] at (7.4,-1.45)
    {\footnotesize next generation: use cost ranking, not policy gradients};
\end{tikzpicture}}%
}

\newcommand{\LostSalesEventFigure}{%
\resizebox{0.98\linewidth}{!}{%
\begin{tikzpicture}[x=1cm,y=1cm]
  \node[invevsource] (sup) at (0,0) {external\\supplier};
  \node[invevstock] (stk) at (5.0,0) {on-hand\\stock $I_t$};
  \node[invevdemand] (dem) at (10.0,0) {demand\\$D_t$};
  \node[invevlost] (lost) at (10.0,-2.7) {lost\\demand};

  \draw[invrequest] (stk.north) to[bend right=36] node[invlabel,above] {$q_t$ request} (sup.north);
  \draw[invmaterial] (sup.south) to[bend right=36] node[invlabel,below] {$q_{t-L}$ arrives} (stk.south);
  \draw[invdemandflow] (stk) -- node[invlabel,above] {serve $\min(I_t,D_t)$} (dem);
  \draw[invlostflow] (dem) -- node[invlabel,right] {unmet} (lost);

  \node[invcost,invblue] at (5.0,-2.7) {holding\\$h(I_t-D_t)^+$};
  \node[invcost,invred,right=2pt of lost.east] {penalty\\$p(D_t-I_t)^+$};

  \begin{scope}[on background layer]
    \node[invbox,fit=(sup)(stk)(dem)(lost),inner sep=10pt] (box) {};
  \end{scope}
  \node[invblue,font=\small\bfseries,above=2pt of box.north]
    {Vanilla lost-sales event sequence};
  \node[invgray,below=4pt of box.south,align=center,font=\footnotesize]
    {arrival $\rightarrow$ observe state $\mathbf S_t$ $\rightarrow$ order request $\rightarrow$ demand realization $\rightarrow$ holding and lost-sales cost};
\end{tikzpicture}}%
}

\newcommand{\FixedCostLostSalesEventFigure}{%
\resizebox{0.98\linewidth}{!}{%
\begin{tikzpicture}[x=1cm,y=1cm]
  \node[invevsource] (sup) at (0,0) {external\\supplier};
  \node[invevstock] (stk) at (5.0,0) {on-hand\\stock $I_t$};
  \node[invevdemand] (dem) at (10.0,0) {demand\\$D_t$};
  \node[invevlost] (lost) at (10.0,-2.7) {lost\\demand};

  \draw[invrequest] (stk.north) to[bend right=36] node[invlabel,above] {$q_t$ request} (sup.north);
  \draw[invmaterial] (sup.south) to[bend right=36] node[invlabel,below] {$q_{t-L}$ arrives} (stk.south);
  \draw[invdemandflow] (stk) -- node[invlabel,above] {serve $\min(I_t,D_t)$} (dem);
  \draw[invlostflow] (dem) -- node[invlabel,right] {unmet} (lost);

  \node[invcost,invblue] at (5.0,-2.7) {holding\\$h(I_t-D_t)^+$};
  \node[invcost,invred,right=2pt of lost.east] {penalty\\$p(D_t-I_t)^+$};
  \node[invcost,invgray,left=2pt of sup.west] {setup\\$K\mathbf{1}\{q_t>0\}$};

  \begin{scope}[on background layer]
    \node[invbox,fit=(sup)(stk)(dem)(lost),inner sep=10pt] (box) {};
  \end{scope}
  \node[invblue,font=\small\bfseries,above=2pt of box.north]
    {Fixed-cost lost-sales event sequence};
  \node[invgray,below=4pt of box.south,align=center,font=\footnotesize]
    {same physical sequence as vanilla lost sales, with a setup charge whenever the requested order is positive};
\end{tikzpicture}}%
}
